\documentclass[10pt]{article}
\usepackage[margin=20mm]{geometry}
\usepackage{amsmath,amssymb,fontspec,unicode-math,xurl,hyperref}
\IfFontExistsTF{DejaVu Serif}{\setmainfont{DejaVu Serif}}{\setmainfont{Latin Modern Roman}}
\setmathfont{latinmodern-math.otf}
\setlength{\parindent}{0pt}\setlength{\parskip}{6pt}\sloppy
\begin{document}
Abstract: The printed core algorithm increases, rather than minimizes, its convex cross-entropy objective on an admissible two-source example with exact Bayes models. This is a Major specification error even after correcting a separate unbound prefactor. The released implementation uses the correct descent sign, so the finding does not establish failure of the reported experiments. The inspected Bayes-mixture identity and uniform approximation proof are valid in their stated setting. A separate static release entry-point defect is Minor. Scientific adjudication is pending.\par
Author: AI-assisted technical review\par
Date:\par
Title: Technical Verification of
MixMin: Finding Data Mixtures via Convex Minimization\par
\# Identity, version, and scope\par
This report concerns Anvith Thudi, Evianne Rovers, Yangjun Ruan, Tristan Thrush, and Chris J. Maddison, MixMin: Finding Data Mixtures via Convex Minimization, ICML 2025, PMLR 267, publisher record `thudi25a` (OpenReview `wpaxYGgp2n`; local identifier 0236/I25-43605). The staged publisher PDF has 17 pages and was retrieved on 25 October 2026 UTC. Page references use printed numbers, equal to one-based PDF indices. The inspected source is `Anvith-Thudi/MixMin`, commit `7212abef3b6e62d2a6164675a46f2ab143d305ca`. Its identity is preserved; it is not independently established as the exact historical experiment revision.\par
The parent read the main paper, both proofs on p. 12, and supplementary result tables and captions on pp. 13–17; relevant released files were inspected statically with auxiliary assistance. Source p. 4 was rendered and visually checked. Own scalar calculations were run, but author code, model inference, training and dataset evaluation were not. The main findings concern mathematical specification and released source, not a reproduced benchmark. A possible separate OpenReview attachment remained inaccessible in source preparation (HTTP 403); embedded appendices were available.\par
\# Verified reduction and limits\par
For unconditional cross-entropy, the Bayes density of a mixture is \(\sum_p\lambda_p f_p\). For conditional cross-entropy or squared loss, the exact conditional Bayes prediction is \[f_\lambda(x)=
 \frac{\sum_p\lambda_p p_p(x) f_p(x)}
      {\sum_p\lambda_p p_p(x)}\] on its support. When all source input marginals agree, this reduces to \(\sum_p\lambda_p f_p(x)\). Theorem 3.1 explicitly imposes the appropriate loss and common-marginal conditions and assumes the relevant Bayes models are in the hypothesis class. Appendix A.1’s substitution is correct in this setting; convexity follows from convex loss composed with an affine mixture.\par
For Lemma 3.2, let \(F_H\) be the finite-class target risk and \(F_*\) the Bayes version. Its uniform approximation and Lipschitz assumptions give \(|F_H(\lambda)-F_*(\lambda)|\le C\epsilon\) for every mixture. If \(\lambda^*\) minimizes \(F_*\) and \(\bar\lambda\) minimizes \(F_H\), then \[F_H(\lambda^*)-F_H(\bar\lambda)
 \le C\epsilon+[F_*(\lambda^*)-F_*(\bar\lambda)]+C\epsilon
 \le 2C\epsilon.\] This verifies the stated argument. Merely enlarging a model class does not by itself establish the uniform approximation assumption or target-risk regularity; the paper states these conditions rather than proving them for the experimental neural networks or bioassays. Source-support extrapolation and empirical Bayes proximity are not certified here.\par
\# Substantive findings\par
Problem 1 (The printed minimization algorithm moves toward the worse source).  Algorithm 1, p. 4, defines \(g=\nabla_\lambda l\) but updates \(\lambda_p\leftarrow\lambda_p e^{+\eta g_p}/\sum_q\lambda_q e^{+\eta g_q}\). For positive step size this has the opposite sign to entropic descent.\par
\#\#\#\# Exact admissible check.\par
Take two distributions on \(\{a,b\}\), \(q_1=(0.9,0.1)\) and \(q_2=(0.1,0.9)\), with their exact Bayes density models. Let the target put all mass on \(a\) and let the hypothesis class include all source mixtures. This is unconditional cross-entropy, so no conditional covariate-shift assumption is needed. A one-sample target dataset containing \(a\) exactly represents its population risk. For \(\lambda=(w,1-w)\), the intended objective is \[l(w)=-\log(0.1+0.8w),\qquad 0\le w\le1.\] Algorithm 1 initializes \(w=1/2\). The unrestricted mixture-coordinate gradient is \(g=(-1.8,-0.2)\), so with the paper’s experimental step size \(\eta=1\) its printed update gives \[w'=\frac{e^{-1.8}}{e^{-1.8}+e^{-0.2}}
 =\frac{1}{1+e^{1.6}}\approx0.1679816149.\] Thus \[l(1/2)=0.6931471806,\qquad l(w')=1.4507889711.\] This is an increase in the exact convex objective, not stochastic noise or a proxy approximation effect. More generally, for every interior \(w\), \[\log\frac{w'}{1-w'}=
 \log\frac{w}{1-w}-\frac{0.8\eta}{0.1+0.8w}.\] Every positive step decreases \(w\), and repeated steps converge to the worse source. The limiting loss is \(-\log0.1\), whereas the minimum is \(-\log0.9\); the excess is \(\log9\).\par
Source visual inspection also confirms that line 3 has an extra \(\lambda_p\) prefactor outside the empirical average, with \(p\) unbound at that step. The counterexample above charitably removes that prefactor to use the intended objective displayed in Section 4. Consequently the sign defect remains even after this separate specification ambiguity is resolved.\par
\#\#\#\# Dependency, counterevidence, and severity.\par
Section 4 identifies Algorithm 1 as the gradient method that solves the paper’s central convex objective. Its failure on exact Bayes sources invalidates that core printed algorithm as a minimizer. The supported severity is Major, bounded to this algorithmic claim; easy repair does not make the demonstrated optimization failure inconsequential.\par
There is decisive counterevidence against extending this finding to experimental failure: , in `update\_with\_stop`, uses `exp(-lr*grad)` and normalizes the weights. The chemistry update also has the negative exponent. Neither Theorem 3.1 nor Lemma 3.2 depends on the erroneous printed sign. The supplied release therefore already embodies the natural sign correction, although its exact relationship to historical runs remains unverified.\par
\#\#\#\# Verified correction.\par
Remove the stray prefactor in line 3 and replace the exponent by \(-\eta g_p\) for the stated loss gradient. The entropic proximal problem \[\arg\min_{\lambda\in\Delta}
 \left\{\eta\langle g,\lambda\rangle+
       \sum_p\lambda_p\log\frac{\lambda_p}{\lambda_p^{\rm old}}\right\}\] has the normalized exponential solution with that negative sign. On the same example it yields \(w'=0.8320183851\) and \(l(w')=0.2670762278\), a strict decrease. This establishes the sign repair and its behavior in the counterexample; it does not certify arbitrary fixed-step convergence for all practical objectives. The independent `c.py`/`calc.json` records the original and corrected calculations.\par
Problem 2 (A released likelihood-generation entry point reads an undefined argument).  The checked `logprobs.py` script cannot reach its intended likelihood computation under its ordinary command-line entry point, even if its imports succeed.\par
In , the parser registers `tokens`, `output\_dir`, `task` and `model\_size`, then calls `parse\_args()` and immediately reads `args.mixture`. No such argument, parser default or prior namespace assignment is present in the complete script. The namespace attribute is therefore absent. This is a Minor static release reproducibility defect; no author program was executed to demonstrate it, and it does not establish that the historical experiments failed. The apparently unused assignment could be removed or an intended parameter explicitly registered, but a complete repaired author run was not performed or authorized.\par
\# Source mapping, empirical scope, and conclusion\par
The language-model objective in the inspected release mixes sequence probabilities: `logprobs.py` multiplies mean token loss by the number of scored causal tokens; `mixmin\_gen.py` applies log-sum-exp across source sequence scores, then length normalization. It is not the unsupported alternative of averaging source next-token probabilities with fixed weights independently at every position. Initial-token conditioning and document-length weighting remain implementation conventions rather than a new theorem guarantee.\par
The chemistry source excludes the target assay when assembling surrogate models and uses the target training split to fit weights; released example weights are proportional to mixture weight divided by source sample count. These static mappings do not establish absence of molecular overlap or verify AP gains. The auxiliary record also notes that long OpenWebMath documents in one branch reuse preceding likelihood values; the measured frequency and historical effect are unresolved. No numerical damage to a reported figure is inferred from that release observation. Supplementary p-values include nonsignificant comparisons, so descriptive improvement is not treated as uniform statistical significance.\par
Overall assessment: Major for the demonstrated core printed algorithm error, with a separate Minor release entry-point defect. The reduction and approximation proof survive the inspected checks. Scientific adjudication remains pending, `human\_reviews` is empty, and neither empirical repair completion nor author acceptance is claimed.\par
9 A. Thudi, E. Rovers, Y. Ruan, T. Thrush, and C. J. Maddison. MixMin: Finding Data Mixtures via Convex Minimization. PMLR 267 (2025). Frozen publisher PDF, 17 pages. https://proceedings.mlr.press/v267/thudi25a.html. MixMin authors. Released source, commit `7212abef3b6e62d2a6164675a46f2ab143d305ca`. https://github.com/Anvith-Thudi/MixMin.\par
\# Evidence and roles\par
Source PDF SHA-256: . Original archives and source mappings remain unchanged. The verification directory contains `inputs.json`, `coverage.json`, `findings.json`, `provenance.json`, independent calculations, source image `s04.png`, and auxiliary `a\_math.json`/`a\_code.json`. Native auxiliaries supplied proof and source checks. The parent independently checked the decisive claims, code and calculations, compiled the exact final TeX and inspected all PDF pages. Actual model visibility, hashes, compiler and visual records are in `provenance.json`. No author code, installations or new datasets were used.\par
\end{document}
